\subsection{Caché LRU de bloques descomprimidos}

La gestión temporal de bloques materializados en memoria se implementó mediante la
clase \texttt{ChunkCache}, cuya función es conservar un conjunto acotado de bloques
descomprimidos recientemente utilizados. Su incorporación evita repetir ciclos completos
de recuperación cuando varias operaciones consecutivas requieren acceder a las mismas
regiones del vector de estado.

La caché se configura según una capacidad determinada y administra la adquisición de
bloques mediante una política LRU. Cuando un bloque es requerido, el backend consulta
primero si ya se encuentra disponible en la caché. Si está presente, se reutiliza
directamente su buffer descomprimido. Si no está presente, el bloque debe recuperarse
desde el almacenamiento comprimido y, en caso de que la caché esté llena, expulsar el
bloque menos recientemente utilizado. Si el bloque expulsado fue modificado desde su
ingreso a la caché, se recompone su versión comprimida antes de liberar el buffer
asociado.

La utilidad de esta implementación no se reduce a mejorar el rendimiento. La caché
establece una forma explícita de administrar la memoria densa temporal del backend y
determina qué parte del estado permanece materializada entre una compuerta y la
siguiente. Por ello, su comportamiento influye directamente en la cantidad de
descompresiones, en la frecuencia de escrituras y en la sobrecarga efectiva del sistema.
